\chapter{Conclusion}
This thesis has developed two applications of complex analysis to integration problems, the inversion of the Laplace transform through contour integration, and the asymptotic approximation of complex integrals by the method of steepest descent (also known as the saddle point approximation). In both cases the aim has been to complete the computations from first principles instead of using a pre-existing table of results, and instead to understand the analytic and geometric structure that underlies and justifies each step in the process.
\\
\\
The section of the inverse Laplace transform began with the derivation of the Bromwich integral from the Fourier inversion formula and then the establishment of the half-plane structure of the region of convergence which were all required for the method to follow. The complex inversion formula was then developed in three cases of increasing complexity. In the simplest case, where the transform has finitely many poles, the inversion would reduce cleanly to a sum of residues once the Bromwich contour is closed and the arc integral is shown to vanish. In the branch point case, the contour must be replaced by a keyhole contour or a suitably designed alternative that renders the integrand single-valued instead of remaining multi-valued, and the inversion produces an integral representation rather than a discrete residue sum, a structural difference that reflects the different nature of the singularity. In the infinitely-many-poles case the essential method remains unchanged, but the contour must be chosen at each stage so that the arc passes between consecutive poles, and the result is an infinite series whose convergence must be verified separately.
\\
\\
The method of steepest descent was developed in its general form, covering saddlepoints of arbitrary order and the distinction between interior and boundary saddlepoints. The central idea is that the integral is controlled entirely by the local geometry of the integrand at the saddlepoint, and that a change of variables converts the oscillatory complex integral into a real Laplace-type integral which was derived carefully and expressed in the general asymptotic formula involving the Gamma function. The standard case was then recovered as a special case, and the method was applied to concrete examples.
\\
\\
What ties both of these topics together is the advantage gained from moving into the complex plane even when the original problem is completely real. In the Laplace inversion, the ability to deform the contour freely when justified by Cauchy's theorem, provided singularities are avoided is what makes the residue calculations possible. In the method of steepest descent, it is the freedom to deform the contour through the saddlepoint that converts the oscillatory integral into a decaying integral. In both cases the key insight of the method comes from understanding the geometry of the complex plane and using it effectively. 
%A recurring theme across both topics is the power that comes from working in the complex plane even when the original problem is entirely real. In the Laplace inversion, the ability to deform contours freely when justified by Cauchy's theorem, provided singularities are avoided is what makes the residue calculation possible. In the steepest descent method, it is precisely the freedom to deform the contour through a saddle point that converts an intractable oscillatory integral into a tractable decaying one. In both cases the key insight is geometric: the structure of the complex plane, rather than brute-force calculation, is what makes the problem tractable. 
\\
\\
Natural directions for further work include the following. A more systematic and rigorous treatment of the convergence of the residue series in the infinitely-many poles case, which would be a direct continuation of section \ref{sec3.1.5}. For the example in this section, the convergence was clear from the exponential decay of the terms, but for transforms with more slowly decaying poles this requires more justification, and the question of whether the series converges uniformly in $t$ or only pointwise is also not addressed. 
A second direction is the extension of the the method of steepest descent method to higher order terms in the asymptotic expansion. 
A third direction would the application of both methods contained in this thesis onto special functions. For example the Airy function, Bessel functions, and the Hankel functions which all have integral representations for which the method of steepest descent applies directly , and working through these would then show how the general asymptotic formula \ref{Eq:steepest_result} produces the results that are usually stated without derivation. 
The Riemann zeta function similarly admits asymptotic analysis through saddlepoint methods which could connect my material contained in this thesis to analytic number theory.
\\
\\
Going beyond the purely mathematical extensions discussed above, the techniques that were developed in this thesis have a wide range of real world applications and connections. The pole based inversion method used in section \ref{Sec3.3} is fundamental in electrical engineering, where the Laplace transform is used to analyse RLC circuits by converting differential equations for voltage and current into algebraic ones in the $s$-domain, with the inverse recovering the time domain response via the exact residue calculations used in this thesis. The same framework underpins control systems in engineering, where the location poles in the complex plane directly determines the stability of the system, connecting back to the half-plane structure showed at the start of Chapter 4. The branch point case is not just a technical complication but can arise naturally in heat and diffusion problems. The same class of branch point singularities can be seen in viscoelasticity, where the stress relaxation of materials produces transforms involving fractional powers of $s$, which is part of what motivated the Jacobs application in section \ref{jacobs}. The method of steepest descent in Chapter 5 is also seen in wave problems, where the field diffracted by an obstacle is expressed as a contour integral with large parameter $k$ playing the exact role of the large parameter throughout the chapter, and the saddlepoints correspond with the points of stationary phase of the wave. The same method also underlies the WKB approximation in quantum mechanics, where approximate solutions to the Schrödinger equation in slowly varying potentials are obtained by applying steepest descent to the path of the integral representation of the propagator, with saddlepoints corresponding to classical particle trajectories. In fluid mechanics, the Kelvin ship wake pattern is derived by applying the saddlepoint method to the dispersion integral for water waves, where the group velocity of the wave packet is determined by the saddlepoint of the phase function. What all of these applications share is the same underlying principle shared with this thesis: that moving into a complex plane, even when the original problem is entirely real reveals a structure that would otherwise remain hidden.